% !TEX root = ../main.tex
%% Chapitre 12 — Espaces vectoriels et applications linéaires
%% Programme officiel MPSI/MP2I (BO spécial n°1 du 11/02/2021), deuxième semestre :
%% A) espaces vectoriels (structure, sous-espaces, Vect, familles, sommes directes) ;
%% B) dimension finie (bases extraite/incomplète, dimension, Grassmann,
%% supplémentaires) ; C) applications linéaires (noyau, image, détermination,
%% endomorphismes, projecteurs, symétries, théorème du rang, formes linéaires et
%% hyperplans) ; D) sous-espaces affines et équations linéaires u(x) = a.
%% Résultats vérifiés par verifications/chapitre_12.py.

\begin{ouverture}[espaces]{12}{Espaces vectoriels et applications linéaires}{Partie 2 · Deuxième semestre}
  {Structure\enspace·\enspace Vect\enspace·\enspace Bases et dimension\enspace·\enspace
   Sommes directes\enspace·\enspace Noyau et image\enspace·\enspace Théorème du rang}
  \objectif{Reconnaître un $\K$-espace vectoriel, un sous-espace, un $\operatorname{Vect}$.}
  \objectif{Manier familles libres, génératrices, bases et coordonnées.}
  \objectif{Calculer des dimensions : base extraite, base incomplète, formule de Grassmann.}
  \objectif{Décomposer un espace en somme directe et exploiter les supplémentaires.}
  \objectif{Étudier une application linéaire : noyau, image, rang, projecteurs et symétries.}
  \objectif{Résoudre géométriquement une équation linéaire $u(x)=a$ : sous-espaces affines.}
  \prerequis{Structures algébriques : groupes, anneaux, corps, morphismes (chapitre 8).}
  \prerequis{Systèmes linéaires et pivot de Gauss (chapitre 9) ; polynômes (chapitre 10).}
  \prerequis{Équations différentielles linéaires (chapitre 4) ; suites récurrentes (chapitre 5).}
\end{ouverture}

%% =====================================================================
%%  L'ESSENTIEL
%% =====================================================================
\essentiel

\essSection{1}{La structure d'espace vectoriel}

\begin{definition}[Espace vectoriel sur $\K$]
Soit $\K=\R$ ou $\Cx$. Un \textbf{$\K$-espace vectoriel} est un ensemble $E$
muni d'une addition qui en fait un groupe commutatif (neutre $0_E$) et d'une
multiplication externe $(\lambda,x)\mapsto\lambda x$ vérifiant, pour tous
$\lambda,\mu\in\K$ et $x,y\in E$ :
\[
\lambda(x+y)=\lambda x+\lambda y,\quad
(\lambda+\mu)x=\lambda x+\mu x,\quad
\lambda(\mu x)=(\lambda\mu)x,\quad
1\cdot x=x .
\]
Les éléments de $E$ sont les \emph{vecteurs}, ceux de $\K$ les
\emph{scalaires}. On a $\lambda x=0_E$ si et seulement si $\lambda=0$ ou
$x=0_E$.
\end{definition}

\begin{exemple}[Le paysage des espaces usuels]
$\K^{n}$ (les $n$-uplets), $\K[X]$ (les polynômes), $\mathcal{M}_{n,p}(\K)$
(les matrices), $\K^{\N}$ (les suites), et plus généralement l'ensemble
$\K^{\Omega}$ des fonctions d'un ensemble $\Omega$ dans $\K$ — ou dans
n'importe quel espace vectoriel. Un produit fini $E_1\times\dots\times E_k$
d'espaces vectoriels en est un, composante par composante. Toute l'algèbre
linéaire s'applique d'un coup à ces mondes en apparence très différents.
\end{exemple}

\begin{definition}[Combinaison linéaire]
Une \textbf{combinaison linéaire} d'une famille finie $(x_1,\dots,x_p)$ de
vecteurs de $E$ est un vecteur $\lambda_1x_1+\dots+\lambda_px_p$ avec
$\lambda_i\in\K$. Pour une famille infinie $(x_i)_{i\in I}$, on impose à la
famille de scalaires $(\lambda_i)_{i\in I}$ d'être \emph{presque nulle}
(nulle sauf pour un nombre fini d'indices), et la somme est alors finie.
\end{definition}

\essSection{2}{Sous-espaces vectoriels et $\operatorname{Vect}$}

\begin{theoreme}[Caractérisation des sous-espaces]
Une partie $F$ de $E$ est un \textbf{sous-espace vectoriel} si et seulement
si $0_E\in F$ et $F$ est stable par combinaison linéaire :
\[
\forall x,y\in F,\ \forall\lambda\in\K,\qquad x+\lambda y\in F .
\]
Un sous-espace est lui-même un espace vectoriel. Toute intersection de
sous-espaces vectoriels est un sous-espace vectoriel.
\end{theoreme}

\begin{exemple}
Le sous-espace nul $\ens{0_E}$ ; une droite vectorielle $\K u$ ($u\neq0$) ;
un plan vectoriel de $\R^{3}$ ; $\K_n[X]$ (degré $\leqslant n$) dans
$\K[X]$ ; l'ensemble des solutions d'un système linéaire \emph{homogène}
dans $\K^{p}$ ; les suites bornées dans $\R^{\N}$. En revanche une réunion
de deux droites distinctes n'est jamais un sous-espace : la somme de deux
vecteurs directeurs sort des deux droites.
\end{exemple}

\begin{definition}[Sous-espace engendré]
Pour $A\subseteq E$, $\operatorname{Vect}(A)$ est l'ensemble des
combinaisons linéaires (finies) de vecteurs de $A$. C'est le \emph{plus
petit} sous-espace vectoriel contenant $A$ : tout sous-espace contenant $A$
contient $\operatorname{Vect}(A)$. On note
$\operatorname{Vect}(x_1,\dots,x_p)$ pour une famille finie.
\end{definition}

\begin{savoirfaire}[Montrer qu'une partie est un sous-espace]
Deux voies, toujours les mêmes : vérifier $0\in F$ puis la stabilité
$x+\lambda y\in F$ — ou, plus rapide quand c'est possible, reconnaître $F$
comme un $\operatorname{Vect}$, une intersection de sous-espaces connus, ou
le noyau d'une application linéaire (section 6).
\end{savoirfaire}

\essSection{3}{Familles libres, génératrices, bases}

\begin{definition}[Libre, génératrice, base]
Une famille $(x_1,\dots,x_p)$ de vecteurs de $E$ est :
\begin{itemize}[label={\color{otNavy}\textbullet},leftmargin=1.4em,itemsep=0.6mm,topsep=0.8mm]
  \item \textbf{génératrice} de $E$ si $\operatorname{Vect}(x_1,\dots,x_p)=E$ ;
  \item \textbf{libre} si $\lambda_1x_1+\dots+\lambda_px_p=0$ entraîne
    $\lambda_1=\dots=\lambda_p=0$ (sinon elle est \emph{liée}) ;
  \item une \textbf{base} de $E$ si elle est libre et génératrice : tout
    $x\in E$ s'écrit alors de façon \emph{unique}
    $x=\lambda_1x_1+\dots+\lambda_px_p$, et les $\lambda_i$ sont les
    \textbf{coordonnées} de $x$ dans la base.
\end{itemize}
\end{definition}

\begin{propriete}[Outils de liberté]
Une famille d'un vecteur est libre si et seulement s'il est non nul ; deux
vecteurs sont liés si et seulement s'ils sont colinéaires. Si
$(x_1,\dots,x_p)$ est libre et $x\notin\operatorname{Vect}(x_1,\dots,x_p)$,
alors $(x_1,\dots,x_p,x)$ est libre. Enfin, toute famille de polynômes non
nuls \emph{à degrés deux à deux distincts} est libre — en particulier les
familles « à degrés échelonnés » $(P_0,\dots,P_n)$ avec $\deg P_k=k$, qui
forment des bases de $\K_n[X]$.
\end{propriete}

\begin{exemple}[Bases canoniques]
$\K^{n}$ : les $e_i=(0,\dots,1,\dots,0)$ ; $\K_n[X]$ : $(1,X,\dots,X^{n})$ ;
$\K[X]$ : $(X^{k})_{k\in\N}$ ; $\mathcal{M}_{n,p}(\K)$ : les matrices
élémentaires $E_{i,j}$. Et $(1,X-1,(X-1)^{2},\dots,(X-1)^{n})$ est une
autre base de $\K_n[X]$, échelonnée en degré : les coordonnées y sont
données par la formule de Taylor (chapitre 10).
\end{exemple}

\essSection{4}{Dimension finie}

\begin{theoreme}[Existence de bases]
$E$ est dit \textbf{de dimension finie} s'il possède une famille génératrice
finie. Alors :
\begin{itemize}[label={\color{otNavy}\textbullet},leftmargin=1.4em,itemsep=0.6mm,topsep=0.8mm]
  \item \emph{base extraite} : de toute famille génératrice finie on peut
    extraire une base ;
  \item \emph{base incomplète} : toute famille libre peut être complétée en
    une base, en piochant dans une famille génératrice donnée.
\end{itemize}
\end{theoreme}

\begin{theoreme}[Dimension]
Dans un espace engendré par $n$ vecteurs, toute famille de $n+1$ vecteurs
est liée. Toutes les bases de $E$ ont donc le même nombre d'éléments, la
\textbf{dimension} $\dim E$. En dimension $n$, une famille de $n$ vecteurs
est une base dès qu'elle est libre \emph{ou} génératrice : la moitié du
travail suffit.
\end{theoreme}

\begin{propriete}[Dimensions à connaître]
$\dim\K^{n}=n$ ; $\dim\K_n[X]=n+1$ ; $\dim\mathcal{M}_{n,p}(\K)=np$ ;
$\dim(E\times F)=\dim E+\dim F$. L'ensemble des solutions de $y'=a(x)y$ est
une droite vectorielle ; celui des solutions de $y''+ay'+by=0$ (coefficients
constants), comme celui des suites vérifiant $u_{n+2}=au_{n+1}+bu_n$, est de
dimension $2$ : deux conditions initiales déterminent tout. Le \textbf{rang}
$\operatorname{rg}(x_1,\dots,x_p)$ est la dimension de
$\operatorname{Vect}(x_1,\dots,x_p)$.
\end{propriete}

\essSection{5}{Sous-espaces, sommes et supplémentaires}

\noindent
\begin{minipage}[c]{0.55\linewidth}
\begin{theoreme}[Sommes et sommes directes]
La \textbf{somme} $F+G=\ens{f+g\mid f\in F,\ g\in G}$ est un sous-espace.
Elle est \textbf{directe}, notée $F\oplus G$, si la décomposition de tout
vecteur de $F+G$ en $f+g$ est \emph{unique}, ce qui équivaut à
$F\cap G=\ens{0}$. On dit que $F$ et $G$ sont \textbf{supplémentaires} dans
$E$ si $E=F\oplus G$ : tout vecteur de $E$ se décompose alors de façon
unique $x=f+g$.
\end{theoreme}
\end{minipage}\hfill
\begin{minipage}[c]{0.41\linewidth}
\centering
\begin{tikzpicture}[scale=0.92,line cap=round]
  \fill[otNavy,opacity=0.06] (-1.7,-0.55) -- (1.15,-0.95) -- (1.7,0.55) -- (-1.15,0.95) -- cycle;
  \draw[otNavy,line width=0.7pt] (-1.7,-0.55) -- (1.15,-0.95) -- (1.7,0.55) -- (-1.15,0.95) -- cycle;
  \draw[otGreenDark,line width=1pt] (-0.42,-1.45) -- (0.72,2.5);
  \draw[otMuted,line width=0.6pt,dash pattern=on 2pt off 1.5pt] (1.05,-0.33) -- (1.45,1.05);
  \draw[otMuted,line width=0.6pt,dash pattern=on 2pt off 1.5pt] (0.4,1.38) -- (1.45,1.05);
  \draw[otNavy,line width=0.8pt,-{Stealth[length=1.9mm]}] (0,0) -- (1.05,-0.33);
  \draw[otGreenDark,line width=0.8pt,-{Stealth[length=1.9mm]}] (0,0) -- (0.4,1.38);
  \draw[otInk,line width=1.1pt,-{Stealth[length=2.2mm]}] (0,0) -- (1.45,1.05);
  \fill[otInk] (0,0) circle (1.3pt);
  \node[font=\scriptsize,text=otInk,anchor=west] at (1.5,1.12) {$x$};
  \node[font=\scriptsize,text=otNavy] at (1.08,-0.6) {$f$};
  \node[font=\scriptsize,text=otGreenDark,anchor=east] at (0.32,1.45) {$g$};
  \node[font=\scriptsize,text=otNavy] at (-1.35,-0.78) {$F$};
  \node[font=\scriptsize,text=otGreenDark,anchor=west] at (0.74,2.35) {$G$};
\end{tikzpicture}\\[0.5mm]
\legende{$E=F\oplus G$ : chaque $x$ se décompose de façon unique en
$f+g$, comme un plan et une droite non incluse dans $\R^{3}$.}
\end{minipage}

\begin{theoreme}[Dimension des sous-espaces ; formule de Grassmann]
Si $F$ est un sous-espace de $E$ de dimension finie, alors
$\dim F\leqslant\dim E$, avec égalité si et seulement si $F=E$. Pour deux
sous-espaces de dimension finie :
\[
\dim(F+G)=\dim F+\dim G-\dim(F\cap G) .
\]
Tout sous-espace $F$ de $E$ possède un supplémentaire, et
\[
E=F\oplus G
\iff
F\cap G=\ens{0}\ \text{et}\ \dim F+\dim G=\dim E
\iff
F+G=E\ \text{et}\ \dim F+\dim G=\dim E .
\]
Concaténer une base de $F$ et une base de $G$ donne alors une base de $E$,
dite \emph{adaptée} à la décomposition.
\end{theoreme}

\begin{savoirfaire}[Montrer que $E=F\oplus G$]
En dimension finie, le plus efficace est presque toujours :
$F\cap G=\ens{0}$ (un seul petit calcul) \emph{plus} le compte des
dimensions $\dim F+\dim G=\dim E$. L'analyse-synthèse — supposer $x=f+g$,
en déduire $f$ et $g$, puis vérifier — reste la méthode universelle, valable
même en dimension infinie (fonctions paires et impaires, exercice 10).
\end{savoirfaire}

\essSection{6}{Applications linéaires : noyau, image, rang}

\begin{definition}[Application linéaire]
$u:E\to F$ est \textbf{linéaire} si $u(x+\lambda y)=u(x)+\lambda u(y)$ pour
tous $x,y\in E$, $\lambda\in\K$. L'ensemble $\mathcal{L}(E,F)$ des
applications linéaires est un espace vectoriel ; la composée de deux
applications linéaires est linéaire, et la composition est bilinéaire. Un
\textbf{isomorphisme} est une application linéaire bijective ; sa réciproque
est alors automatiquement linéaire.
\end{definition}

\begin{theoreme}[Noyau et image]
Pour $u\in\mathcal{L}(E,F)$ :
$\operatorname{Ker}u=\ens{x\in E\mid u(x)=0}$ est un sous-espace de $E$,
$\operatorname{Im}u=u(E)$ un sous-espace de $F$ — plus généralement, image
directe et image réciproque d'un sous-espace sont des sous-espaces. Et :
\[
u\ \text{injective}\iff\operatorname{Ker}u=\ens{0},
\qquad
u\ \text{surjective}\iff\operatorname{Im}u=F .
\]
Si $(x_i)_{i\in I}$ engendre $E$, alors
$\operatorname{Im}u=\operatorname{Vect}\bigl(u(x_i)\bigr)_{i\in I}$. Le
\textbf{rang} de $u$ est $\operatorname{rg}u=\dim\operatorname{Im}u$ ; on a
$\operatorname{rg}(v\circ u)\leqslant\min(\operatorname{rg}u,\operatorname{rg}v)$,
et le rang est invariant par composition par un isomorphisme.
\end{theoreme}

\begin{theoreme}[Détermination par une base]
Si $(e_i)_{i\in I}$ est une base de $E$ et $(f_i)_{i\in I}$ une famille
quelconque de $F$, il existe une \emph{unique} $u\in\mathcal{L}(E,F)$ telle
que $u(e_i)=f_i$ pour tout $i$ : une application linéaire est entièrement
décidée par l'image d'une base. De même, si $E=E_1\oplus E_2$, une
application linéaire se définit en la prescrivant sur $E_1$ et sur $E_2$.
Conséquences : deux espaces de dimension finie sont isomorphes si et
seulement s'ils ont la même dimension, et
$\dim\mathcal{L}(E,F)=\dim E\times\dim F$.
\end{theoreme}

\essSection{7}{Endomorphismes : projecteurs et symétries}

\noindent
\begin{minipage}[c]{0.55\linewidth}
\begin{theoreme}[Projecteurs et symétries]
$\bigl(\mathcal{L}(E),+,\circ\bigr)$ est un anneau (non commutatif dès que
$\dim E\geqslant2$) ; ses inversibles forment le \emph{groupe linéaire}
$\operatorname{GL}(E)$. Si $E=F\oplus G$ et $x=f+g$ :
\begin{itemize}[label={\color{otNavy}\textbullet},leftmargin=1.4em,itemsep=0.6mm,topsep=0.8mm]
  \item le \textbf{projecteur} sur $F$ parallèlement à $G$ est $p(x)=f$ ;
    il vérifie $p^{2}=p$, $\operatorname{Im}p=F=\operatorname{Ker}(\mathrm{id}-p)$
    et $\operatorname{Ker}p=G$ ;
  \item la \textbf{symétrie} par rapport à $F$ parallèlement à $G$ est
    $s(x)=f-g=2p(x)-x$ ; elle vérifie $s^{2}=\mathrm{id}$.
\end{itemize}
Réciproquement, tout $p$ vérifiant $p^{2}=p$ est le projecteur sur
$\operatorname{Im}p$ parallèlement à $\operatorname{Ker}p$ (et
$E=\operatorname{Ker}p\oplus\operatorname{Im}p$) ; tout $s$ vérifiant
$s^{2}=\mathrm{id}$ est la symétrie associée à $p=\frac12(s+\mathrm{id})$.
\end{theoreme}
\end{minipage}\hfill
\begin{minipage}[c]{0.41\linewidth}
\centering
\begin{tikzpicture}[scale=1.0,line cap=round]
  \draw[otNavy,line width=1pt] (-1.5,-0.5) -- (1.8,0.6);
  \draw[otGreenDark,line width=0.9pt] (-0.55,-1.5) -- (0.55,1.5);
  \fill[otInk] (0,0) circle (1.2pt);
  \draw[otInk,line width=1pt,-{Stealth[length=2mm]}] (0,0) -- (1.35,1.15);
  \draw[otMuted,line width=0.6pt,dash pattern=on 2pt off 1.5pt] (1.35,1.15) -- (1.05,0.35);
  \draw[otNavy,line width=0.9pt,-{Stealth[length=2mm]}] (0,0) -- (1.05,0.35);
  \draw[otMuted,line width=0.6pt,dash pattern=on 2pt off 1.5pt] (1.35,1.15) -- (0.75,-0.45);
  \draw[otGreenDark,line width=0.9pt,-{Stealth[length=2mm]}] (0,0) -- (0.75,-0.45);
  \node[font=\scriptsize,text=otInk,anchor=south west] at (1.3,1.1) {$x$};
  \node[font=\scriptsize,text=otNavy,anchor=north west] at (1.05,0.3) {$p(x)$};
  \node[font=\scriptsize,text=otGreenDark,anchor=north west] at (0.72,-0.5) {$s(x)$};
  \node[font=\scriptsize,text=otNavy,anchor=west] at (1.82,0.6) {$F$};
  \node[font=\scriptsize,text=otGreenDark,anchor=south] at (0.55,1.52) {$G$};
\end{tikzpicture}\\[0.5mm]
\legende{Projection sur $F$ parallèlement à $G$, et symétrie :
$s(x)=2p(x)-x$. Se faire \emph{toujours} ce dessin.}
\end{minipage}

\essSection{8}{Théorème du rang, hyperplans, sous-espaces affines}

\begin{theoreme}[Théorème du rang]
Soit $u\in\mathcal{L}(E,F)$ avec $\dim E=n$ finie. Si $S$ est un
supplémentaire de $\operatorname{Ker}u$ dans $E$, alors $u$ induit un
\emph{isomorphisme} de $S$ sur $\operatorname{Im}u$ (forme géométrique).
En particulier :
\[
\dim E=\dim\operatorname{Ker}u+\operatorname{rg}u .
\]
Si $\dim E=\dim F$ finies, alors $u$ injective $\iff$ surjective $\iff$
bijective — et un endomorphisme inversible à gauche \emph{ou} à droite est
inversible.
\end{theoreme}

\begin{theoreme}[Formes linéaires et hyperplans]
Une \textbf{forme linéaire} sur $E$ est une application linéaire
$\varphi:E\to\K$ ; dans une base, $\varphi(x)=a_1x_1+\dots+a_nx_n$. Un
\textbf{hyperplan} est le noyau d'une forme linéaire non nulle. Si $H$ est
un hyperplan et $D$ une droite non incluse dans $H$, alors $E=H\oplus D$ ;
en dimension $n$, les hyperplans sont exactement les sous-espaces de
dimension $n-1$. Deux équations $\varphi(x)=0$ et $\psi(x)=0$ définissent le
même hyperplan si et seulement si $\varphi$ et $\psi$ sont proportionnelles.
Enfin, l'intersection de $m$ hyperplans est de dimension au moins $n-m$, et
tout sous-espace de dimension $n-m$ est intersection de $m$ hyperplans :
c'est la lecture géométrique des systèmes linéaires.
\end{theoreme}

\begin{definition}[Sous-espaces affines]
Un \textbf{sous-espace affine} de $E$ est une partie $\mathcal{A}=x_0+F$ où
$F$ est un sous-espace vectoriel, la \emph{direction} de $\mathcal{A}$
(l'écriture $B=A+\vec u$ se lit $\overrightarrow{AB}=\vec u$). Deux
sous-espaces affines de même direction sont égaux ou disjoints ; une
intersection non vide de sous-espaces affines est un sous-espace affine.
\end{definition}

\begin{theoreme}[Structure des solutions de $u(x)=a$]
Soit $u\in\mathcal{L}(E,F)$ et $a\in F$. L'ensemble des solutions de
$u(x)=a$ est soit vide, soit le sous-espace affine
$x_0+\operatorname{Ker}u$, où $x_0$ est une solution particulière :
\emph{solution générale = solution particulière + solution homogène}.
Systèmes linéaires, équations différentielles linéaires, suites
arithmético-géométriques, interpolation : un seul et même énoncé.
\end{theoreme}

%% =====================================================================
%%  EXERCICES
%% =====================================================================
\exercices
\rubrique{Applications}

\begin{exercice}[1]{Sous-espace ou pas ?}
Les parties suivantes sont-elles des sous-espaces vectoriels ?
\begin{questions}
  \item $\ens{(x,y)\in\R^{2}\mid x+2y=0}$ et
    $\ens{(x,y)\in\R^{2}\mid x+y=1}$ ;
  \item $\ens{(x,y)\in\R^{2}\mid xy=0}$ ;
  \item $\ens{f\in\R^{\R}\mid f(1)=0}$ et $\ens{f\in\R^{\R}\mid f(1)=1}$ ;
  \item $\ens{P\in\R[X]\mid\deg P=3}$ ;
  \item l'ensemble des suites réelles bornées.
\end{questions}
\end{exercice}

\begin{exercice}[1]{Dans un Vect, ou pas ?}
Dans $\R^{3}$, on pose $u=(1,1,1)$ et $v=(0,1,2)$.
\begin{questions}
  \item Le vecteur $(1,2,3)$ appartient-il à $\operatorname{Vect}(u,v)$ ?
    Et $(1,0,0)$ ?
  \item Décrire $\operatorname{Vect}(u,v)$ par une équation cartésienne.
  \item À quelle condition sur $c$ le vecteur $(1,2,c)$ est-il dans
    $\operatorname{Vect}(u,v)$ ?
\end{questions}
\end{exercice}

\begin{exercice}[1]{Libres ou liées ?}
Étudier la liberté des familles suivantes :
\begin{questions}
  \item $\bigl((1,0,1),(1,1,0),(0,1,1)\bigr)$ dans $\R^{3}$ ;
  \item $\bigl((1,2,3),(2,4,6)\bigr)$ dans $\R^{3}$ ;
  \item $\bigl((1,1,0),(0,1,1),(1,0,-1)\bigr)$ dans $\R^{3}$
    (donner une relation de liaison) ;
  \item $(1,X,X^{2},\dots,X^{n})$ dans $\R[X]$, puis $(\cos,\sin)$ dans
    $\R^{\R}$.
\end{questions}
\end{exercice}

\begin{exercice}[1]{Une base échelonnée de $\R_3[X]$}
\begin{questions}
  \item Montrer que $\mathcal{B}=\bigl(1,\,X-1,\,(X-1)^{2},\,(X-1)^{3}\bigr)$
    est une base de $\R_3[X]$.
  \item Déterminer les coordonnées de $X^{3}$ dans $\mathcal{B}$.
  \item Que donne, en général, la formule de Taylor en $1$ ?
\end{questions}
\end{exercice}

\begin{exercice}[1]{Noyau et image d'une application linéaire}
Soit $u:\R^{3}\to\R^{2}$ définie par $u(x,y,z)=(x+y-z,\ 2x-y)$.
\begin{questions}
  \item Vérifier que $u$ est linéaire.
  \item Déterminer une base de $\operatorname{Ker}u$.
  \item Déterminer $\operatorname{Im}u$ et $\operatorname{rg}u$. Le
    théorème du rang est-il confirmé ?
\end{questions}
\end{exercice}

\begin{exercice}[1]{Linéaire ou pas ?}
Les applications suivantes sont-elles linéaires ?
\begin{questions}
  \item $f(x,y)=(x+y,\ x-y)$ et $g(x,y)=(x+1,\ y)$ sur $\R^{2}$ ;
  \item $h(x,y)=xy$ de $\R^{2}$ dans $\R$ ;
  \item la dérivation $P\mapsto P'$ sur $\R[X]$, et la translation
    $P\mapsto P(X+1)$ ;
  \item l'évaluation $\varphi:f\mapsto f(2)$ de $\R^{\R}$ dans $\R$.
\end{questions}
\end{exercice}

\begin{exercice}[1]{Un projecteur à la loupe}
Soit $p:\R^{2}\to\R^{2}$ définie par
$p(x,y)=\Bigl(\dfrac{x+y}{2},\ \dfrac{x+y}{2}\Bigr)$.
\begin{questions}
  \item Vérifier que $p$ est linéaire et que $p\circ p=p$.
  \item Déterminer $\operatorname{Im}p$ et $\operatorname{Ker}p$ ; de quel
    projecteur s'agit-il ? Faire un dessin.
  \item On pose $s=2p-\mathrm{id}$. Calculer $s(x,y)$, vérifier
    $s\circ s=\mathrm{id}$ et interpréter géométriquement.
\end{questions}
\end{exercice}

\begin{exercice}[1]{Rangs et dimensions}
\begin{questions}
  \item Déterminer $\operatorname{rg}\bigl((1,1,0),(0,1,1),(1,0,-1)\bigr)$
    et $\operatorname{rg}\bigl((1,2),(2,4),(3,6)\bigr)$.
  \item Quelle est la dimension de $\ens{(x,y,z)\in\R^{3}\mid x+y+z=0}$ ?
    En donner une base.
  \item Montrer que $(1+X,\,1-X,\,X^{2})$ est une base de $\R_2[X]$.
\end{questions}
\end{exercice}

\rubrique{Entraînement}

\begin{exercice}[2]{Des espaces de solutions}
\begin{questions}
  \item Rappeler pourquoi l'ensemble des solutions de $y''+y=0$ est un
    espace vectoriel, et vérifier que $(\cos,\sin)$ en est une famille
    libre. Admettant que cet espace est de dimension $2$, qu'en déduit-on ?
  \item Soit $E$ l'ensemble des suites réelles vérifiant
    $u_{n+2}=u_{n+1}+u_n$. Montrer que
    $u\mapsto(u_0,u_1)$ est un isomorphisme de $E$ sur $\R^{2}$ ; en
    déduire $\dim E=2$.
  \item Montrer que les suites géométriques $(\varphi^{n})$ et
    $(\psi^{n})$, où $\varphi=\frac{1+\sqrt5}{2}$ et
    $\psi=\frac{1-\sqrt5}{2}$, forment une base de $E$, et retrouver la
    formule de Binet pour la suite de Fibonacci (chapitre 5).
\end{questions}
\end{exercice}

\begin{exercice}[2]{Supplémentaires : deux grands classiques}
\begin{questions}
  \item Dans $\R^{3}$, on pose $F=\ens{(x,y,z)\mid x+y+z=0}$ et
    $G=\operatorname{Vect}\bigl((1,1,1)\bigr)$. Montrer que
    $\R^{3}=F\oplus G$ et décomposer $(1,2,3)$.
  \item Dans $\R^{\R}$, montrer que toute fonction s'écrit de façon unique
    comme somme d'une fonction paire et d'une fonction impaire. Que donne
    la décomposition pour $\exp$ ?
\end{questions}
\end{exercice}

\begin{exercice}[2]{La formule de Grassmann en action}
\begin{questions}
  \item Montrer que deux plans vectoriels distincts de $\R^{3}$ se coupent
    toujours selon une droite.
  \item Dans $\R^{3}$, expliciter cette droite pour
    $F=\operatorname{Vect}\bigl((1,0,0),(0,1,0)\bigr)$ et
    $G=\operatorname{Vect}\bigl((0,0,1),(1,1,1)\bigr)$.
  \item Deux sous-espaces de dimension $3$ de $\R^{5}$ peuvent-ils être
    d'intersection nulle ?
\end{questions}
\end{exercice}

\begin{exercice}[2]{Base incomplète, base extraite}
\begin{questions}
  \item Compléter la famille libre $\bigl((1,1,0),(0,1,1)\bigr)$ en une
    base de $\R^{3}$ à l'aide d'un vecteur de la base canonique.
  \item Extraire une base de $\R^{2}$ de la famille génératrice
    $\bigl((1,0),(2,0),(0,1),(1,1)\bigr)$.
  \item Déterminer un supplémentaire de
    $\operatorname{Vect}\bigl(1+X+X^{2},\,X\bigr)$ dans $\R_2[X]$.
\end{questions}
\end{exercice}

\begin{exercice}[2]{Détermination d'une application linéaire}
Soit $u$ l'unique endomorphisme de $\R^{2}$ tel que $u(1,1)=(2,0)$ et
$u(1,-1)=(0,4)$.
\begin{questions}
  \item Pourquoi $u$ existe-t-il et est-il unique ?
  \item Exprimer $u(x,y)$ pour tout $(x,y)\in\R^{2}$, et calculer $u(1,0)$.
  \item Montrer que $u$ est un automorphisme de $\R^{2}$.
\end{questions}
\end{exercice}

\begin{exercice}[2]{Le théorème du rang en pratique}
Soit $u:\R^{4}\to\R^{3}$ dont la matrice dans les bases canoniques est
$A=\begin{pmatrix}1&2&0&1\\0&1&1&1\\1&3&1&2\end{pmatrix}$.
\begin{questions}
  \item Déterminer $\operatorname{rg}u$.
  \item Donner une base de $\operatorname{Ker}u$ et une base de
    $\operatorname{Im}u$, et vérifier le théorème du rang.
  \item L'application $u$ est-elle injective ? surjective ?
\end{questions}
\end{exercice}

\begin{exercice}[2]{Projecteurs et symétries : les caractérisations}
Soit $p\in\mathcal{L}(E)$ un projecteur ($p^{2}=p$).
\begin{questions}
  \item Montrer que $\mathrm{id}-p$ est un projecteur et que
    $\operatorname{Ker}(\mathrm{id}-p)=\operatorname{Im}p$.
  \item Montrer que $E=\operatorname{Ker}p\oplus\operatorname{Im}p$,
    en écrivant $x=\bigl(x-p(x)\bigr)+p(x)$.
  \item Soit $s(x,y)=(y,x)$ sur $\R^{2}$. Vérifier $s^{2}=\mathrm{id}$,
    déterminer les deux droites $\operatorname{Ker}(s-\mathrm{id})$ et
    $\operatorname{Ker}(s+\mathrm{id})$, et reconnaître $s$.
\end{questions}
\end{exercice}

\begin{exercice}[2]{Hyperplans de $\R^{3}$}
Soit $\varphi(x,y,z)=x-2y+z$ et $H=\operatorname{Ker}\varphi$.
\begin{questions}
  \item Justifier que $H$ est un hyperplan et en donner une base.
  \item Trouver une droite $D$ telle que $\R^{3}=H\oplus D$.
  \item Montrer que $-2x+4y-2z=0$ définit le même hyperplan. Deux équations
    non proportionnelles peuvent-elles définir le même hyperplan ?
\end{questions}
\end{exercice}

\begin{exercice}[2]{L'interpolation, vue par l'algèbre linéaire}
Soit $\Phi:\R_2[X]\to\R^{3}$, $P\mapsto\bigl(P(0),P(1),P(2)\bigr)$.
\begin{questions}
  \item Montrer que $\Phi$ est linéaire et injective, puis que c'est un
    isomorphisme.
  \item En déduire, sans aucun calcul, l'existence et l'unicité du polynôme
    de degré au plus $2$ prenant des valeurs prescrites en $0,1,2$
    (chapitre 10).
  \item Expliciter les trois polynômes $L_0,L_1,L_2$ dont les images par
    $\Phi$ forment la base canonique de $\R^{3}$, et vérifier
    $L_0+L_1+L_2=1$.
\end{questions}
\end{exercice}

\begin{exercice}[2]{Solution particulière + solution homogène}
\begin{questions}
  \item Décrire l'ensemble des solutions du système
    $\begin{cases}x-y=1\\x+y+z=3\end{cases}$ comme sous-espace affine de
    $\R^{3}$ : point particulier et direction.
  \item Résoudre $u_{n+1}=2u_n+3$ avec $u_0=1$ par la même méthode :
    suite constante particulière, puis solutions homogènes (chapitre 5).
  \item Même lecture pour $y'=y+1$ : donner la solution générale.
\end{questions}
\end{exercice}

\rubrique{Approfondissement}

\begin{exercice}[3]{Familles libres de fonctions}
\begin{questions}
  \item Soient $a_1<a_2<\dots<a_n$ des réels. Montrer que la famille des
    fonctions $x\mapsto\mathrm{e}^{a_ix}$ est libre dans $\R^{\R}$.
    \emph{Diviser une relation de liaison par $\mathrm{e}^{a_nx}$ et faire
    $x\to+\infty$, puis itérer.}
  \item Montrer que la famille $\bigl(x\mapsto|x-a|\bigr)_{a\in\R}$ est
    libre. \emph{Penser à la dérivabilité en chaque point $a$.}
\end{questions}
\end{exercice}

\begin{exercice}[3]{Endomorphismes nilpotents : premiers pas}
Soit $E$ de dimension $n$ et $u\in\mathcal{L}(E)$ \emph{nilpotent} :
$u^{p}=0$ et $u^{p-1}\neq0$ pour un certain $p\geqslant1$ (l'\emph{indice}).
\begin{questions}
  \item Soit $x$ tel que $u^{p-1}(x)\neq0$. Montrer que
    $\bigl(x,u(x),\dots,u^{p-1}(x)\bigr)$ est libre. En déduire
    $p\leqslant n$.
  \item Vérifier que la dérivation est nilpotente sur $\R_n[X]$, d'indice
    $n+1$.
  \item Montrer que $\mathrm{id}-u$ est inversible, d'inverse
    $\mathrm{id}+u+u^{2}+\dots+u^{p-1}$ (chapitres 8 et 9).
\end{questions}
\end{exercice}

\begin{exercice}[3]{Noyaux itérés}
Soit $u\in\mathcal{L}(E)$, $E$ de dimension finie $n$.
\begin{questions}
  \item Montrer que la suite des noyaux
    $\operatorname{Ker}u^{0}\subseteq\operatorname{Ker}u\subseteq
    \operatorname{Ker}u^{2}\subseteq\dots$ est croissante.
  \item Montrer que si $\operatorname{Ker}u^{k}=\operatorname{Ker}u^{k+1}$
    pour un indice $k$, alors $\operatorname{Ker}u^{m}=\operatorname{Ker}u^{k}$
    pour tout $m\geqslant k$.
  \item En déduire que la suite stationne au plus tard au rang $n$.
  \item Expliciter la suite des noyaux pour la dérivation sur $\R_3[X]$.
\end{questions}
\end{exercice}

\begin{exercice}[3]{Deux projecteurs qui commutent}
Soient $p,q$ deux projecteurs de $E$ tels que $p\circ q=q\circ p$.
\begin{questions}
  \item Montrer que $p\circ q$ est un projecteur.
  \item Montrer que
    $\operatorname{Im}(p\circ q)=\operatorname{Im}p\cap\operatorname{Im}q$.
    \emph{Utiliser la caractérisation $y\in\operatorname{Im}p\iff p(y)=y$.}
  \item Montrer que
    $\operatorname{Ker}(p\circ q)=\operatorname{Ker}p+\operatorname{Ker}q$.
    \emph{Pour l'inclusion directe, écrire $x=\bigl(x-q(x)\bigr)+q(x)$.}
\end{questions}
\end{exercice}

\begin{exercice}[3]{Les hyperplans de polynômes}
Dans $E=\R_3[X]$ (de dimension $4$), on pose, pour $a\in\R$,
$H_a=\ens{P\in E\mid P(a)=0}$.
\begin{questions}
  \item Montrer que $H_a$ est un hyperplan de $E$ et que
    $H_a=(X-a)\,\R_2[X]$.
  \item Montrer que $H_0\cap H_1\cap H_2$ est une droite vectorielle, dont
    on donnera un générateur.
  \item Plus généralement, que vaut
    $\dim(H_{a_0}\cap\dots\cap H_{a_m})$ pour $a_0,\dots,a_m$ deux à deux
    distincts ($m\leqslant3$) ? Commenter avec l'énoncé « l'intersection
    de $m$ hyperplans est de dimension au moins $n-m$ ».
\end{questions}
\end{exercice}

\begin{exercice}[3]{Quand $E=\operatorname{Ker}u\oplus\operatorname{Im}u$}
Soit $E$ de dimension finie et $u\in\mathcal{L}(E)$. Montrer l'équivalence
des quatre assertions :
\begin{questions}
  \item $E=\operatorname{Ker}u\oplus\operatorname{Im}u$ ;
  \item $\operatorname{Ker}u^{2}=\operatorname{Ker}u$ ;
  \item $\operatorname{Im}u^{2}=\operatorname{Im}u$ ;
  \item $\operatorname{rg}u^{2}=\operatorname{rg}u$.
\end{questions}
Vérifier que tout projecteur les satisfait, mais que la dérivation sur
$\R_2[X]$ ne satisfait aucune d'elles.
\end{exercice}

% !TEX root = ../../main.tex
%% Défi du chapitre 12 — L'opérateur de différence et la formule de Newton
%% Résultats vérifiés par verifications/defis_11_15.py.

\rubrique{Défi}

\begin{defi}{L'opérateur de différence et la formule de Newton}{1 h 30}
Soit $n\in\N^{*}$ et $E=\R_n[X]$. On définit l'opérateur de différence
$\Delta:P\mapsto P(X+1)-P(X)$, version « discrète » de la dérivation, et
les polynômes
\[
H_0=1,\qquad H_k=\frac{X(X-1)\cdots(X-k+1)}{k!}\quad(k\geqslant1),
\]
qui vérifient $H_k(m)=\dbinom mk$ pour tout entier $m\geqslant0$.

\partieExo{A}{Un endomorphisme nilpotent}
\begin{questions}
  \item Montrer que $\Delta$ est un endomorphisme de $E$, et que si
    $\deg P=d\geqslant1$, alors $\deg\Delta P=d-1$ : préciser le
    coefficient dominant de $\Delta P$.
  \item Déterminer $\operatorname{Ker}\Delta$ et $\operatorname{Im}\Delta$
    (théorème du rang). Montrer que $\Delta^{n+1}=0$ mais
    $\Delta^{n}\neq0$.
  \item Écrire la matrice de $\Delta$ dans la base canonique de $\R_3[X]$.
\end{questions}

\partieExo{B}{La base binomiale et la formule de Newton}
\begin{questions}[resume]
  \item Montrer que $(H_0,\dots,H_n)$ est une base de $E$, et que
    $\Delta H_k=H_{k-1}$ pour $k\geqslant1$ : c'est la formule de Pascal
    pour les polynômes. Écrire la matrice de $\Delta$ dans cette base.
  \item En déduire la \emph{formule de Newton} : pour tout $P\in E$,
    \[
    P=\sum_{k=0}^{n}\bigl(\Delta^{k}P\bigr)(0)\,H_k .
    \]
    Décomposer ainsi $X^{3}$, en lisant les $\bigl(\Delta^{k}P\bigr)(0)$
    sur la table des différences ci-dessous.
  \item Montrer que $P$ prend des valeurs entières en tout entier relatif
    si et seulement si ses coordonnées dans la base $(H_k)$ sont entières.
    Exemple : $\frac{X(X+1)}{2}$, à coefficients non entiers, mais à
    valeurs entières.
\end{questions}

\medskip
\begin{center}
\begin{tikzpicture}[x=1.05cm,y=0.72cm]
  \foreach \v [count=\i from 0] in {0,1,8,27,64,125}{%
    \node[font=\small,text=otInk] at (\i,0) {$\v$};}
  \foreach \v [count=\i from 0] in {1,7,19,37,61}{%
    \node[font=\small,text=otNavy] at (\i+0.5,-1) {$\v$};}
  \foreach \v [count=\i from 0] in {6,12,18,24}{%
    \node[font=\small,text=otNavy] at (\i+1,-2) {$\v$};}
  \foreach \v [count=\i from 0] in {6,6,6}{%
    \node[font=\small,text=otNavy] at (\i+1.5,-3) {$\v$};}
  \foreach \v [count=\i from 0] in {0,0}{%
    \node[font=\small,text=otMuted] at (\i+2,-4) {$\v$};}
  \foreach \xx/\yy in {0/0,0.5/-1,1/-2,1.5/-3}{%
    \draw[otGreenDark,line width=0.8pt,rounded corners=2pt] (\xx-0.27,\yy-0.3) rectangle (\xx+0.27,\yy+0.3);}
  \foreach \l/\yy in {$P(k)=k^{3}$/0,$\Delta P$/-1,$\Delta^{2}P$/-2,$\Delta^{3}P$/-3,$\Delta^{4}P$/-4}{%
    \node[font=\scriptsize,text=otMuted,anchor=east] at (-0.6,\yy) {\l};}
\end{tikzpicture}\\[0.5mm]
\legende{La table des différences de $k\mapsto k^{3}$ : chaque ligne est
formée des différences de la précédente ; la diagonale encadrée donne
$\bigl(\Delta^{k}P\bigr)(0)=0,1,6,6$.}
\end{center}

\partieExo{C}{Sommation discrète}
\begin{questions}[resume]
  \item Montrer que si $\Delta Q=P$, alors
    $\sum_{k=0}^{m-1}P(k)=Q(m)-Q(0)$ pour tout $m\in\N^{*}$ :
    $Q$ joue le rôle d'une « primitive discrète » de $P$.
  \item Trouver $Q$ tel que $\Delta Q=X^{2}$, et en déduire
    $\sum_{k=0}^{m-1}k^{2}$. Faire de même pour
    $\sum_{k=0}^{m-1}k^{3}$, et retrouver le théorème de
    Nicomaque (défi du chapitre 2).
\end{questions}
\end{defi}

